\documentclass[11pt]{article}
\usepackage[T1]{fontenc}
\usepackage{textcomp}
\usepackage[margin=0.75in]{geometry}
\usepackage{amsmath,amssymb,amsthm}
\usepackage{array,booktabs}
\usepackage{hyperref}
\setlength{\parindent}{0pt}
\newtheorem{theorem}{Theorem}
\theoremstyle{definition}
\newtheorem{definition}{Definition}
\title{Flow Proof Artifact: Nat-order}
\author{Generated by \texttt{flow doc proof}}
\date{Flow Proof Book}
\begin{document}
\maketitle
Order axioms and basic lemmas on natural numbers.\\[0.5em]
\textbf{Source.} peano — https://en.wikipedia.org/wiki/Total\_order\\[0.5em]
\bigskip
\noindent{\large \textbf{Derived fact 1.} \textit{Every natural is less than or equal to itself} \hfill the natural numbers \cdot order \cdot less-or-equal is reflexive}\par
\smallskip\noindent\textit{Coordinate: the natural numbers \cdot order \cdot less-or-equal is reflexive}\par
\smallskip\noindent\textit{Source: peano}\par
\medskip
\noindent\textbf{Goal.} Every natural is less than or equal to itself.  n ≤ n\par
\noindent$\displaystyle \forall n \in \mathbb{N}\quad n \le n$\par
\medskip
\noindent
\begin{tabular}{@{} >{\raggedright\arraybackslash}p{0.06\textwidth} >{\raggedright\arraybackslash}p{0.47\textwidth} >{\raggedleft\arraybackslash}p{0.41\textwidth} @{}}
\textbf{\#} & \textbf{Proof} & \textbf{Mathematics} \\
\midrule
\textbf{1.} & We prove that less-or-equal is reflexive for order on the natural numbers. & \\[0.45em]
\textbf{2.} & We can deduce that n is at most n. Hence proven. & $\displaystyle n \le n$ \\[0.45em]
\bottomrule
\end{tabular}
\label{thm:Nat-order-less_or_equal_is_reflexive}
\par\medskip\hrule
\bigskip
\noindent{\large \textbf{Derived fact 2.} \textit{If a ≤ b and b ≤ a, then a = b} \hfill the natural numbers \cdot order \cdot less-or-equal is antisymmetric}\par
\smallskip\noindent\textit{Coordinate: the natural numbers \cdot order \cdot less-or-equal is antisymmetric}\par
\smallskip\noindent\textit{Source: peano}\par
\smallskip\noindent\textit{Built on: left cancellation holds, for addition on the natural numbers}\par
\medskip
\noindent\textbf{Goal.} If a ≤ b and b ≤ a, then a = b\par
\noindent$\displaystyle \forall a \in \mathbb{N} \forall b \in \mathbb{N}\quad a = b$\par
\medskip
\noindent
\begin{tabular}{@{} >{\raggedright\arraybackslash}p{0.06\textwidth} >{\raggedright\arraybackslash}p{0.47\textwidth} >{\raggedleft\arraybackslash}p{0.41\textwidth} @{}}
\textbf{\#} & \textbf{Proof} & \textbf{Mathematics} \\
\midrule
\textbf{1.} & We prove that less-or-equal is antisymmetric for order on the natural numbers. & \\[0.45em]
\textbf{2.} & We split into exhaustive cases — the claim must hold in each one. & \\[0.45em]
\textbf{3.} & Case 1 (see step 2): suppose a <= b. & \\[0.45em]
\textbf{4.} & Case 2 (see step 2): suppose b <= a. & \\[0.45em]
\textbf{5.} & We invoke the derived fact governing addition on the natural numbers: left cancellation holds, for addition on the natural numbers (instantiated for a, 0, 0). & \\[0.45em]
\textbf{6.} & From step 4 and step 5, this implies a equals b. Together with the other cases (step 3 and step 4), the goal is discharged. Hence proven. & $\displaystyle a = b$ \\[0.45em]
\bottomrule
\end{tabular}
\label{thm:Nat-order-less_or_equal_is_antisymmetric}
\par\medskip\hrule
\bigskip
\noindent{\large \textbf{Derived fact 3.} \textit{If a ≤ b and b ≤ c, then a ≤ c} \hfill the natural numbers \cdot order \cdot less-or-equal is transitive}\par
\smallskip\noindent\textit{Coordinate: the natural numbers \cdot order \cdot less-or-equal is transitive}\par
\smallskip\noindent\textit{Source: peano}\par
\medskip
\noindent\textbf{Goal.} If a ≤ b and b ≤ c, then a ≤ c\par
\noindent$\displaystyle \forall a \in \mathbb{N} \forall b \in \mathbb{N} \forall c \in \mathbb{N}\quad a \le c$\par
\medskip
\noindent
\begin{tabular}{@{} >{\raggedright\arraybackslash}p{0.06\textwidth} >{\raggedright\arraybackslash}p{0.47\textwidth} >{\raggedleft\arraybackslash}p{0.41\textwidth} @{}}
\textbf{\#} & \textbf{Proof} & \textbf{Mathematics} \\
\midrule
\textbf{1.} & We prove that less-or-equal is transitive for order on the natural numbers. & \\[0.45em]
\textbf{2.} & We split into exhaustive cases — the claim must hold in each one. & \\[0.45em]
\textbf{3.} & Case 1 (see step 2): suppose a <= b. & \\[0.45em]
\textbf{4.} & Case 2 (see step 2): suppose b <= c. & \\[0.45em]
\textbf{5.} & From step 4, this implies a is at most c. Together with the other cases (step 3 and step 4), the goal is discharged. Hence proven. & $\displaystyle a \le c$ \\[0.45em]
\bottomrule
\end{tabular}
\label{thm:Nat-order-less_or_equal_is_transitive}
\par\medskip\hrule
\bigskip
\noindent{\large \textbf{Derived fact 4.} \textit{Every natural is strictly less than its successor} \hfill the natural numbers \cdot order \cdot every number is below its successor}\par
\smallskip\noindent\textit{Coordinate: the natural numbers \cdot order \cdot every number is below its successor}\par
\smallskip\noindent\textit{Source: peano}\par
\smallskip\noindent\textit{Built on: adding zero on the left does not change the number}\par
\medskip
\noindent\textbf{Goal.} Every natural is strictly less than its successor.  n < succ(n)\par
\noindent$\displaystyle \forall n \in \mathbb{N}\quad n < \mathrm{succ}(n)$\par
\medskip
\noindent
\begin{tabular}{@{} >{\raggedright\arraybackslash}p{0.06\textwidth} >{\raggedright\arraybackslash}p{0.47\textwidth} >{\raggedleft\arraybackslash}p{0.41\textwidth} @{}}
\textbf{\#} & \textbf{Proof} & \textbf{Mathematics} \\
\midrule
\textbf{1.} & We prove that every number is below its successor for order on the natural numbers. & \\[0.45em]
\textbf{2.} & We invoke the definitional clause governing addition on the natural numbers: adding zero on the left does not change the number (instantiated for n). & $\displaystyle 0 + n = n$ \\[0.45em]
\textbf{3.} & From step 2, this implies n < the successor of n. Hence proven. & $\displaystyle n < \mathrm{succ}(n)$ \\[0.45em]
\bottomrule
\end{tabular}
\label{thm:Nat-order-every_number_is_below_its_successor}
\par\medskip\hrule
\end{document}
